% =====================================================================
%  Methods  --  Nature Methods places this after the references.
% =====================================================================
\section*{Methods}

\subsection*{Data and provenance}

All CAID3 numbers derive from the official reference FASTAs
(\url{https://caid.idpcentral.org/assets/sections/challenge/static/references/3/})
and the official prediction archive (\code{predictions.zip}, 76{,}480{,}806
bytes, 117 \code{.caid} files, aggregate md5
\code{cdef7ae6dc3dd6ec8055aacdb14a54a7}). Per-file checksums and per-target
composition counts are in Supplementary Table~S12. Composition is cross-checked
against CAID's dataset API, which serves those counts separately from the files
themselves; the evaluator refuses any reference whose composition does not
match. Of the 117 prediction files, 115 produce a score on any given reference,
which is the denominator of every rank reported here; the capacity tables use
117, since a method that entered is a method the benchmark was asked to place.

The dataset used is ``CAID3 v3'', the challenge as scored. The API also serves
an entry named plainly ``CAID3'' (185 proteins), a different and smaller set that
does not reproduce the published leaderboard. Three prediction files
(\code{FoldUnfold}, \code{NeProc-binding}, \code{NeProc-disorder}) leave the
score field blank on declined residues; these parse to NaN, are counted, and are
never silently dropped.

The check that licenses every other number is that each challenge's published
leader is recomputed through our own pipeline. All five reproduce to the
published precision: PUNCH2 0.9552 against 0.955, ESMDisPred-2PDB 0.8855 against
0.885, DisoFLAG-PB 0.7760 against 0.776, bindEmbed21IDR-rawGeneral 0.6407
against 0.641, IPA-AF2-Linker 0.8985 against 0.897. If a leader stops
reproducing, nothing computed afterwards is trusted until it is explained.

\subsection*{Annotation error rate}

$\eps$ is estimated two independent ways from MobiDB, neither of which uses
CAID's own references (which are not independent annotations of each other and
return exactly zero disagreements).

\emph{Context-dependent residues.} MobiDB's
\code{derived-missing\_residues\_context\_dependent-th\_90} annotation marks
residues missing in some deposited structures of a protein and observed in
others. Pooled over the 160 CAID3 Disorder-PDB targets with a usable record:
4{,}886 of 61{,}013 evidenced residues, $\eps = 0.0801$.

\emph{Structure pairs.} For each protein, every pair of its deposited structures
(\code{derived-missing\_residues-mobi-\{PDBID\}\_\{CHAIN\}}), restricted to
residues both structures cover, gives a direct disagreement rate. Over the 159
targets with $\ge 2$ structures and 2{,}746 structure pairs:
$\eps_{\text{label}} = 0.0651$ pooled and 0.0342 median protein. The same
enumeration gives $\eps_{\text{pair}} = 0.0100$ pooled and 0.0006 median, on the
denominator \code{card\_comparablePairs} defines --- the pairs both structures
order, $2(|T \cap L|\cdot|(T \cup L)^{c}| + du)$, and no larger set. The same run
records the balance hypothesis rather than assuming it: 32 of 2{,}746 structure
pairs satisfy both hypotheses of \code{nuPair\_le\_two\_eps\_sq}, and the bound
holds on all 32.

Per-structure coverage is approximated by the span of that structure's annotated
region, because MobiDB publishes per-structure \emph{missing} regions and not
per-structure coverage; residues outside a structure's span are treated as
undetermined rather than observed, which is the conservative direction, since
counting them as observed would manufacture disagreement from absence. Both
rates are worst-case budgets over structure pairs, not probability models,
matching the noise model the capacity theorem uses.

\subsection*{Formal development}

All theorems cited are Lean~4, \code{sorry}-free, and depend only on
\code{propext}, \code{Classical.choice} and \code{Quot.sound}. The development
covers:

\begin{itemize}\itemsep1pt \parskip0pt
\item the AUC decomposition and its invariance
(\code{auc\_pooled\_decomp};\ \code{auc\_within\_strictMono\_invariant};\
\code{auc\_pooled\_shift\_diff};\ \code{inversion\_requires\_between\_gap});
\item label noise and certification
(\code{LabelNoise.ranking\_certified};\ \code{RobustCertificate});
\item benchmark capacity in three score models
(\code{card\_le\_benchCapacity};\ \code{benchCapacity\_attained};\
\code{benchCapacity\_noise\_only};\ \code{unresolvable\_pair};\
\code{over\_capacity\_has\_close\_pair});
\item pairwise scoring
(\code{discordant\_eq\_flip\_product};\ \code{discordant\_le\_noise\_sq};\
\code{card\_comparablePairs};\ \code{noise\_orderKey};\
\code{pairwise\_capacity\_bound};\ \code{nuPair\_le\_two\_eps\_sq});
\item calibration
(\code{Calibration.risk\_decomposition};\ \code{validity\_is\_free});
\item multiple testing under dependence
(\code{wgt\_decomp};\ \code{selfConsistent\_fdr\_le\_harmonic};\
\code{bh\_fdr\_eq\_harmonic};\ \code{harmonic\_factor\_sharp};\
\code{fdp\_lift};\ \code{region\_screen\_fdr\_le};\
\code{harmonic\_gain\_log});
\item complexity (\code{biasThreshold\_hard};\
\code{crossed\_bound\_loose\_unbounded}).
\end{itemize}

A full inventory is Supplementary Note~S6.

Where the development supplies a worked example the analysis code reproduces it
exactly --- \code{Example.inversion} returns within 1 / pooled 2/3 against within
2/3 / pooled 5/6 --- which is the strongest available check that the Python
computes the statistic the theorems describe.

Two statements are used ahead of their formalisation and flagged as such:
\code{auc\_target\_strictMono\_invariant}, the per-target form of the AUC
invariance, which is what the protocol's step 3 needs (the published theorem is
stated for the pair-weighted mean, and the invariance holds target by target
before any averaging); and the validity of conformal risk control, cited rather
than formalised.

\subsection*{The pairwise protocol}

For a reference with targets $t$ and per-residue labels in $\{0,1,-\}$:

\begin{enumerate}\itemsep1pt \parskip0pt
\item \emph{Eligibility.} A method is scored iff it supplies a finite
prediction for every evaluated residue of every target, at the reference's
length. Declining targets raises a within-protein score, so coverage is a gate,
not a covariate. Ineligible methods are listed as \emph{not scored}, never as
last.
\item \emph{Per-target statistic.} For each target carrying both classes among
evaluated residues, the Mann--Whitney AUC over that target's residues alone.
Single-class targets contribute no ordered pair and are skipped, with the count
reported.
\item \emph{Method score.} The unweighted mean of the per-target AUCs. One
protein, one vote; pair-weighting lets a few long chains carry the number.
\item \emph{Ranking.} By that mean, descending.
\item \emph{Separation.} Two methods are reported as ordered iff a paired test
over targets rejects equality --- Wilcoxon signed-rank on the per-target
differences, Holm-corrected within the benchmark. Otherwise they are reported as
tied, at the same rank.
\item \emph{Capacity.} The number of methods the reference can order is
$\lceil 1/2\eps_{\text{pair}} \rceil$, with $\eps_{\text{pair}}$ estimated from
repeat determinations of the same protein. Methods beyond that are reported as an
unresolved group, never as a rank.
\end{enumerate}

The unweighted mean of step 3 differs from the pair-weighted
$\mathrm{AUC}_{\text{within}}$ that the decomposition identity requires. Both are
calibration-invariant, since the invariance holds target by target before any
averaging; the unweighted form is what a paired per-target design naturally uses
and is not dominated by long chains. A separation count is not a chain: the
capacity of step 6 bounds a certified totally ordered family, whereas $k$
separations against one leader are $k$ statements about the pairs
$\{\text{leader}, X\}$, which need not compose into a chain of $k{+}1$.

\subsection*{Model}

\emph{Backbone.} ESM-2 650M\cite{esm2}, frozen. A learned scalar mixture over
its 33 hidden layers (33 parameters) replaces fine-tuning. The design follows a
measurement: in the small-data regime here ($\sim$1M evidenced residues from
2{,}340 proteins) a 69.9M-parameter LoRA configuration reached pooled AUC 0.7454
while a gradient-boosted tree on hand-crafted features reached 0.7804 and
AlphaFold pLDDT alone reached 0.7906. The bottleneck was capacity/data fit, not
the backbone.

\emph{Head.} Four residual dilated-convolution blocks, hidden width 256,
dilations $(1,4,16,32)$ --- a 213-residue convolutional receptive field for the
same parameter count as the default $(1,2,4,8)$. The \emph{dependency span} is
global, not 213: the head's GroupNorm pools statistics over the whole length
axis, so every output position already depends on every input position. This is
verified directly (\code{lite\_head.measured\_dependency\_span}: on a
401-residue input, perturbing residue 0 moves the logit at residue 400).

\emph{Structure channels.} Relative solvent accessibility, pLDDT (scaled),
contact density, an availability flag, and the local gradient of solvent
accessibility, encoded into a small learned block. Missing structure is explicit
rather than imputed: an absent AlphaFold entry is not ``buried and confident'',
it is no information. Post-hoc fusion of the structural signal failed
($0.9581 \rightarrow 0.9554$); feeding structure as an input lets the trunk
learn when to trust it.

\emph{Read-outs.} One $1\times1$ convolution per task ($\sim$257 parameters
each) over the shared trunk, deliberately linear: anything deeper would let a
task with 15{,}683 positive residues (Linker) grow private capacity, which is
the failure mode the architecture exists to avoid. The shared trunk carries
disorder's 336{,}014 positive residues into the small tasks.

\emph{Windowing.} Training and inference use a 1{,}022-residue window at stride
511 with a raised-cosine taper on the overlap. Every window of a chain is
assigned to the same side of every split, keyed on the parent sequence.

\emph{Variants.} \code{mt\_windowed} (no protein bias, no private trunk; the
pre-registered model of record), \code{mt\_pbias} (adds a per-protein bias
term), and \code{DisorderNet-Ensemble} (equal-weight rank fusion of the
checkpoints sharing the fixed validation holdout, with membership decided on the
holdout and never on CAID3; both fixed in advance).

\subsection*{Leak control}

Splits are homology-clustered with BLASTp, clustered once over the union of
proteins so a protein appearing in two tasks cannot land in different folds for
each. Benchmark targets and their $\ge 40\%$-identity homologues are removed
from training. A fixed 5\% validation holdout is selected by a salted hash of the
parent sequence --- membership is a property of the sequence, not of the run ---
so that two runs with different leak filters remain comparable; the salt is a
module constant with no command-line override, and holdout membership is written
into every checkpoint. The homology filter refuses to run if BLAST is
unavailable rather than holding out nothing, because a validation number
computed with a paralogue in training is worse than none: it looks like a
measurement.

The temporal holdout is built from protein polymer entities first released on or
after 2026-08-11, the day after the later of the two training caches (DisProt
2026-08-08, MobiDB 2026-08-10). Labels are RCSB's own
\code{UNOBSERVED\_RESIDUE\_XYZ} annotation, not a re-derivation. X-ray and
cryo-EM only; NMR is excluded, since every SEQRES residue has coordinates in
every model.

\subsection*{Statistics}

Every p-value belongs to a family fixed in advance and Holm--Bonferroni\cite{holm}
is applied within it --- Holm rather than Bonferroni because it is uniformly more
powerful and assumes no independence, which these comparisons badly lack, sharing
both our predictions and their targets. The primary CAID3 family is fixed in the
pre-registration and encoded in the evaluator so it cannot be redefined after the
numbers land; everything outside it is secondary, corrected within its own
family, and labelled exploratory.

Paired comparisons resample \emph{proteins}, not residues, since residues within
a chain are heavily correlated and a residue-level interval is far too narrow.
All tests are two-sided; 10{,}000 resamples for confirmatory runs. Bootstrap
p-values use both tail boundaries and the Davison--Hinkley $+1$
correction\cite{davison}, so a perfect null returns 1.0 rather than 0 and $p$ is
never reported below $1/(B{+}1)$. Where a figure or table gives $p$, the test,
the family size and the correction are stated with it; exact $n$ (targets,
chains, structure pairs or methods) is given in every legend.

Coverage is a gate, not a covariate. Measured against a fully-covering yardstick,
declined targets are harder by $+0.035$ to $+0.125$ AUC and have median length
1{,}214--1{,}684 against 344 across all CAID3 targets (Supplementary Table~S9).
We predict every target on every benchmark and the evaluator aborts rather than
scoring a subset. Both rank fields are always reported --- rank among all scored
entrants (CAID's own accounting, which embeds the advantage) and rank among
full-coverage entrants (the equal-footing field). Reporting only the first
understates us; reporting only the second is cherry-picking. Neither alone is the
truth.

\subsection*{Conformal analysis}

Split conformal prediction\cite{vovk} and conformal risk control\cite{crc} as
implemented in \code{colab/conformal.py} (34 tests). Scores are rank-transformed
within each method, because conformal needs a consistent score and not a
probability. Calibration and evaluation are on disjoint halves of the targets,
median over 12 protein splits, with the interquartile range recorded. The
implementation refuses by name to report a guarantee at $\alpha < 1/(n{+}1)$,
which is arithmetic about $n$ rather than a fact about any method.

A per-residue promise is \emph{not} valid here, and the shortfall is measured
rather than hidden: residues within a chain are not exchangeable, so the split
is by chain, and at a 90\% target the realised per-residue coverage on the
temporal holdout is 0.876. It is quoted as a measurement of how far the
assumption is strained, not as a guarantee.

\subsection*{Screening analysis}

Two computations. \code{region\_screen.py} counts, for each CAID3 and CAID2
reference, the evaluated residues and the maximal same-label runs, and reports
both harmonic numbers, their difference, the proved floor $\log b - 1$, and the
ratio of the two testing thresholds. A run is broken by any unevaluated residue,
so a region interrupted by missing evidence counts as two candidates --- the
conservative direction, since it states more hypotheses rather than fewer.
Harmonic numbers are summed from the small end.

\code{conformal\_screen.py} runs the screen itself on CAID3 Disorder-PDB (304
usable chains, 98{,}938 evaluated residues). Candidates are consecutive blocks
of $b$ residues, the partition \code{stdBlocks} describes; a block's score is
the mean predicted disorder over its evaluated residues and a block is truly
disordered when the majority of those residues are, so that the truth and the
report are both unions of whole blocks, which is the condition \code{fdp\_lift}
needs. P-values are split conformal against the known-ordered blocks of a
calibration set of 80\% of chains,
$p = (1 + \#\{\text{calibration blocks with score} \ge s\})/(n_{\text{cal}}+1)$,
which is superuniform under exchangeability of the ordered blocks --- the only
assumption \code{selfConsistent\_fdr\_le\_harmonic} makes. It is also why the
correction must survive arbitrary dependence: every candidate is compared against
the \emph{same} calibration set, so the p-values are dependent by construction.
Four procedures are compared at the same guaranteed residue-level rate over 50
random splits (region-level BY at $\alpha/H_M$; residue-level BY at
$\alpha/H_n$; region-level uncorrected BH; region-level Bonferroni), and
\code{fdp\_lift} is checked rather than assumed by printing the region-level and
residue-level false discovery proportions of the same report side by side. The
screen is a demonstration on a public benchmark; no biological discovery is
claimed.

\subsection*{Reproduction}

Figures are regenerated from the analysis outputs by \code{paper/make\_figures.py}
and supplementary tables by \code{paper/make\_supplementary.py}; each panel reads
the JSON its job wrote, and the two panels that quote a table parse it from the
committed result document rather than restating it, so a figure cannot drift from
the number it illustrates. Job identifiers, batch scripts and raw outputs for
every analysis are listed in Supplementary Table~S13.

\subsection*{Reporting summary}

Further information on research design is available in the Nature Portfolio
Reporting Summary linked to this article.
